% =============================================================
% quiz5min.tex — Шаблон пятиминутки по математике
% Компилятор: XeLaTeX (два прохода для корректных crop-линий)
%
% Раскладка страницы:
%   ┌──────────────┬──────────────┐
%   │  Вариант 1   │  Вариант 2   │
%   ├ · · · · · · ·+· · · · · · ·┤  ← линия реза
%   │  Вариант 1   │  Вариант 2   │
%   └──────────────┴──────────────┘
% =============================================================

\documentclass[12pt]{article}

% Путь к стилю — при компиляции из корня latex-math-sheets/
\usepackage{styles/mathworksheet}

% ============================================================
% ПАРАМЕТРЫ КОНКРЕТНОЙ ПЯТИМИНУТКИ
% Меняйте только эти строки при создании новой работы
% ============================================================
\newcommand{\quiztopic}{Сложение и вычитание дробей}
% ============================================================

\begin{document}

% ============================================================
% ЗАДАНИЯ — определяются как макросы, чтобы не дублировать
% В1 и В2 — разные числа, одинаковый тип заданий
% ============================================================

%% ----- ВАРИАНТ 1 -------------------------------------------
\newcommand{\tasksetA}{%
  %% Блок 1: вычислить (3 примера)
  \calcrowc{1}{\tfrac{3}{7} + \tfrac{2}{7}}
  \calcrowc{2}{\tfrac{5}{6} - \tfrac{1}{4}}
  \calcrowc{3}{\tfrac{2}{3} + \tfrac{3}{8}}
  \tasksep
  %% Блок 2: упростить
  \taskrow{4}{Вычислите:\enspace
    $\dfrac{4}{9} + \dfrac{5}{9} - \dfrac{2}{3}$}
  \vspace{4mm}%
  %% Блок 3: текстовая задача
  \taskrowans{5}{Из бочки вылили \mf{\tfrac{3}{8}} объёма,
    затем ещё \mf{\tfrac{1}{4}}. Какую часть вылили?}
}

%% ----- ВАРИАНТ 2 -------------------------------------------
\newcommand{\tasksetB}{%
  \calcrowc{1}{\tfrac{4}{9} + \tfrac{1}{9}}
  \calcrowc{2}{\tfrac{7}{8} - \tfrac{1}{3}}
  \calcrowc{3}{\tfrac{3}{4} + \tfrac{5}{12}}
  \tasksep
  \taskrow{4}{Вычислите:\enspace
    $\dfrac{7}{10} - \dfrac{2}{10} + \dfrac{3}{5}$}
  \vspace{4mm}%
  \taskrowans{5}{Турист прошёл \mf{\tfrac{2}{5}} пути утром
    и \mf{\tfrac{3}{10}} вечером. Какую часть пути прошёл?}
}

% ============================================================
% СЕТКА 2×2 — НЕ ИЗМЕНЯЙТЕ СТРУКТУРУ
% ============================================================
\begin{cardgrid}
  %% Верхний ряд
  \quizcard{\quiztopic}{1}{\tasksetA}
  &
  \quizcard{\quiztopic}{2}{\tasksetB}
  \\[\cardgap]
  %% Нижний ряд (дубликат для второй пары учеников)
  \quizcard{\quiztopic}{1}{\tasksetA}
  &
  \quizcard{\quiztopic}{2}{\tasksetB}
\end{cardgrid}

\end{document}

% ============================================================
% СПРАВОЧНИК КОМАНД (для добавления своих заданий)
% ============================================================
%
%  \calcrowc{N}{выражение}     — вычислительное задание,
%                                textstyle (рекомендуется)
%  \calcrow{N}{выражение}      — то же, displaystyle (крупнее)
%  \taskrow{N}{текст}          — задание с текстом
%  \taskrowans{N}{текст}       — задание + строка для ответа
%  \tasksep                    — лёгкий разделитель блоков
%  \ansline[ширина]            — отдельная линия ответа
%  \mf{формула}                — строчная формула (textstyle)
%  \cardfoot                   — класс + дата в конце карточки
%
% СОВЕТЫ:
%  - Для дробей в строчных заданиях используйте \tfrac{}{}
%  - Для уравнений: \taskrow{N}{Решите уравнение: $3x-7=11$}
%  - 4-5 заданий — оптимум; при 6 заданиях проверьте что
%    карточка не переполнена
%
% КОМПИЛЯЦИЯ:
%  xelatex templates/quiz5min.tex
%  xelatex templates/quiz5min.tex   (второй проход для crop-линий)
% ============================================================
